\chapter{Results}
\label{ch:results}


% Below are the aggregated results from the 10 player qualitative studies that got conducted while writing this thesis.

% I identified 40 themes during the qualitative tests that can be found in the appendix 
% Below are the aggregated findings, meta-themes. 

% While analysing the interview materials I identified 4 themes per participant amounting to 40 themes. From those I managed to find 5 meta-themes -- 
% overarching narratives that can be synthesised from the sum of all of the sessions. They're presented below. \todo{add the documents and table}


% As discussed in Chapter 3 all of the particiapnt

Below are the aggregated results from the user testing described in Chapter~\ref{ch:method}. 
Based on the method described in that chapter, five main categories were formulated.
For the list of all themes that contributed to them, please refer to \ref{tab:responses}.




\section{A - Prior experiences carry-over effect}

\textit{Supporting themes from the analysis: \#1A, \#2A, \#3A, \#10D}

One of the first observations from all the themes was that players used their prior experiences, which had an impact on their interaction with the game. The first three participants who played the keyboard version first transferred that fresh experience into the way they interacted with the voice version of the game, as noted by participant \#1, "because I was guided [in my word choice] by that choice of options” .

Both versions of the game differed substantially when it comes to the user interface. 
The currently available actions that were written in simple language seemed to have been naturally used by the players who were exposed to the text-driven version of the game prior to playing the voice experience. If a player saw an option "Pick up crowbar", they were likely to use the exact same words and language when voicing out their commands. 
Despite switching the order of testing, players still imported past experiences with video games, voice assistants and other input modalities. 




\section{B - Perception of control}

\textit{Supporting themes from the analysis: \#1C, \#2D, \#4B, \#9C, \#10C}


The voice input system had a major impact on the feelings of agency and control. 
It expanded the perceived affordances of the game, allowing the participants to issue commands that are beyond what was supported in the game logic.
That feeling of expanded control clashed with the reality in which the game couldn't fulfil all of the players' requests, given its technical limitations. The effect of that was twofold:

For some participants, the mismatch between their wishes and results created feelings of frustration and anxiety. 
When the game failed to register their commands, the participants reported the sensation of "false sense of control", in which the initial perception of control was undermined, and the trust in the game itself was lowered. Participant \#3 noted, "And I think it gives you the illusion of choice, because you don't know what you can do and what you cannot do," supporting this claim.

Conversely, for other participants, despite registering the feeling of not being fully in charge, the voice input still made them feel more control and freedom compared to the keyboard version. 2nd participant directly supported this claim by saying, "I mean, it [the voice version] wasn't overwhelming for me. It was more like a freedom that I could think of another solution, could try something else."

The fact that the voice input system, unlike the traditional navigation, does not clearly communicate its boundaries resulted in players seeking to find the limits of the system. The emotional impact of this varied between players, ranging from excitement to anxiety and uncertainty. 




\section{C - The relationship between command language and the game experience}

\textit{Supporting themes from the analysis: \#3B, \#5A, \#6A, \#7B}

The participants had different feelings about the game based on the language that they chose to use to interact with it. When initially instructed on how to play, they weren't given a rigid template to follow. Rather than that, they were encouraged to try and experiment with their commands. 

This resulted in participants talking with the game in a varying manner, falling into a few separate groups:

\begin{enumerate}
    \item Cooperative approach -- Some participants treated the game as a co-partner. 
    Participant \#2 directly commented on that dynamic by saying, "When you talk to the computer, you feel like you have a conversation with the game".
    This coincided with them feeling more in control compared to the permission-seeking approach, and more resilient to technical errors as they approached the game like a shared exploratory experience. 
    
    % \begin{quote}
    %     "When you talk to the computer you feel like you have a conversation with the game, so it's also pretty interesting." 
    %      \flushright --- Participant \#2   
    % \end{quote}


    \item First person narration -- Few of the participants expressed their voice commands as direct wishes articulated in first person, akin to playing a tabletop role-playing game. This language can be observed with participant \#6's word choice while playing: "I'm going inside the clothing store, but I'm trying to step inside the steps that I saw on the dusty floor so that I'm covering my tracks". For them, the direct language seems to have resulted in a higher perceived feeling of immersion by the participants compared to the permission seeking group. 
    
    \item Permission seeking -- A large group of participants interacted with the game by posing questions aimed at the game. For instance, participant \#10 would default to this style of communication when feeling uncertain: "Can I use a crowbar?". This created a dynamic that placed the game as an authority that needs to be asked for permission, similar to a work superior or a teacher at school, not a partner.
    This difference from other groups translated to participants feeling less in control compared to the keyboard.
    For some participants, this system was interlaced with the commanding language, gravitating between them during the experience.

    \item Commanding the system -- This style of communication was based on direct commands, similar to commanding a subordinate, as demonstrated by what participant\#7 said while playing: "Look around for any food". This style of communication resulted in a higher feeling of control and confidence compared to seeking permission, similar to the first-person and cooperative interaction styles.
    
    % \begin{quote}
    %     "Look around for any food"
    %     \flushright --- Participant \#7
    % \end{quote}

\end{enumerate}

The reasons why participants chose their interaction style fell into two categories. 
For some participants, it was automatic - they didn't think about it and just naturally picked a style of interaction. 
Other participants consciously chose their interaction style, basing it on their previous experiences with voice assistants. They intentionally chose simple command language since that's what worked best in their experience. 



%  Their choice seemed to have been based on their experiences connected to games, voice assistants and notions about voice control. 

% \begin{quotation}
%     "let's go here, let's go there, I want to do this, I want to do that, that was obvious to me, I mean, I feel like we are playing together, me and the computer, and not like I'm asking him to take me somewhere, it's like we are together in this"
%     \flushright --- Participant \#5 
% \end{quotation}

% These small linguistic differences translated to the players feeling vastly different about the voice modality in general and experiencing the game differently. 
% Those who were asking the game for permission eg. "Can we go here?" felt less in control and more frustrated compared to people who treated the game as a partner "Let's go here". 
% The people who spoke using first person looked to be more immersed into the experience, while those speaking as if the player was a 3rd person character looked to be more present in reality rather than the game. 


% Players all adopted their own language when talking with the computer and nearly always it stemed from their existing relationship with either computers, technology, voice assistants or some other notions that come from speakng to someone for permision. The language that they chose was either commanding or cooperative. This in turn shaped their emotions while playing the game. This type of unpredictable behaviour is absent in the keyboard driven experience. 


\section{D - Emotional impact of technical limitations}

\textit{Supporting themes from the analysis: \#2B, \#4A, \#6C, \#8A, \#10A}


% 4A - playing in public, using english as a non-primary language
% 10A non primary language 

% 6C - understanding gaps 

% 1,2,3 VUI bugs with Lounge or staris 

% 8A - delay 

The voice input generated a range of potential negative emotions in players. The feelings of frustration and annoyance were brought up by multiple participants as a result of them having ot use the VUI. The major source of irritation was the fact that the game would sometimes struggle to understand the participant even after repeated attempts. This, combined with a lack of a clear feedback system on failed requests, left many with a feeling of anger or distrust towards the system.

Another major issue was the higher cognitive load that was required from participants to play the voice-driven version compared to the keyboard navigation. Since the voice-driven version didn't expose the player actions as a list, some players reported feeling overburdened by having to read the full locations' description texts and being forced to think out all of their actions. 
Participant \#6 described his experience as follows: 
\begin{quote}
    Well, thinking what to say takes more effort than pressing a key, obviously.
    I mean, maybe it's not obvious, but...
    It takes more effort because you don't have your actions laid out in front of you and you have an open field, creative open field as of what you can really do.
\end{quote}


Additionally, the processing delays, taking about 5--7 seconds on each request, made some participants feel frustrated and like they were lacking control over the game. For instance, participant \#8 noted that they “forget the things [they] want to do” and, at another point during playing the game, “had something to do but… lost [their] train of thought a little bit”.

On top of that, a few participants noted that having to issue voice commands in English was a much higher linguistic challenge for them. 
This is clearly in contrast to the keyboard version, which was easier to understand language-wise, as it only required reading.
As a result, those participants felt uncertain and anxious.

\section{E - Impact of using voice on player immersion}

% 300 budget
% +
% 1B 2C 5D deeper in experience - reading all text
% 8B deeper experience  
% 2C pausing and having to read it all 
% 5C irrational preference


% ?
% 9B deeper experinece but not prefferd 
% 10B trade off immersion but anxierty and cognitive load

% -
% 3C internalised immersion
% 4C and Kuba speed as immersion gradient - breaking 
% 7A recognition anxiety as supressant 
% 7D cognitive offloading not having to think when navigating with keybard 



\textit{Supporting themes from the analysis: \#1B, \#2C, \#3C, \#4C \#5C, \#5D, \#7A, \#7D, \#8B, \#9B, \#10B}


Using voice-driven navigation had a varying effect on players immersion.

One unintentional reason why this happened was the fact that, in order to navigate in the game, the players were forced to read the entirety of location descriptions. This is more of a side-effect of the voice modality -- forcing the participants to be more present. This observation was clearly noted by the participant who noted that the voice version was more immersive \#2 "because you had to immerse more in the text and more focus on where you are, what you're doing." 

% \begin{quote}
%     "because you had to immerse more in the text and more focus on where you are, what you're doing." 
%     \flushright --- Participant \#2 
% \end{quote}

Despite the fact that when asked for their preference, some participants chose the keyboard version, many still attributed higher immersion to the voice version.

% The keyboard version (due to ease of navigation and responsiveness) they still felt more immersed in the voice version. This showcases that participants would willingly choose an experience that offers quicker and more reliable input over a more immersive one.


Other participants noted that using voice made them feel less immersed in the game compared to the keyboard version. 
Some players described immersion as something close to a flow state - a moment where they "become one" with the game.
They noted that using the voice input broke that state in a way. For some of the few second delay was what broke their trance.
For someone else, the heightened cognitive load made it difficult to maintain immersion, as noted by participant \#7: "reading a whole paragraph is and then trying to figure out what commands you should use by voice is way harder than just going through the list of what you can do."
In another case, the anxiety that the player wouldn't be recognised by the system resulted in a feeling of stress so high that immersion couldn't be maintained. Another participant claimed that the mere act of speaking aloud broke their immersion, since that experience was no longer taking place just in their head in an internalised way. 



% \begin{quote}
%     "reading a whole paragraph is and then trying to figure out what you what commands should you use yeah by voice is way harder than just going through the list of what you can do "
%      \flushright --- Participant \#7 
% \end{quote}




% \section{F - Accessibility context}

% Some players mentioned how this could be a good experience for people who can't use keyboard navigation. 


% \todo{Move to Conclusion for future research}
% \section{G - Communication and learning side benefits}

% one participant mentioned how this could ease someone into speaking in their target language. Another person mentioned how the game could be used to open up the player in general even if they're fluent in their target language.
